% !TEX root = ../main.tex
\begin{ouverture}[spectre]{2}{Méthodes physiques d'analyse}{Partie 1 · Constitution et transformations de la matière}
{Beer-Lambert\enspace·\enspace Kohlrausch\enspace·\enspace Gaz parfait\enspace·\enspace Spectroscopies IR et UV-visible}
\objectif{Exploiter la loi de Beer-Lambert pour déterminer une concentration.}
\objectif{Mesurer une conductivité et exploiter la loi de Kohlrausch ou une droite d'étalonnage.}
\objectif{Déterminer une quantité de matière de gaz avec l'équation d'état du gaz parfait.}
\objectif{Identifier un groupe caractéristique ou une espèce chimique à l'aide d'un spectre IR ou UV-visible.}
\objectif{Citer les domaines de validité de ces relations.}
\prerequis{Absorbance, spectre d'absorption et couleur d'une solution.}
\prerequis{Concentration en quantité de matière, dilution, volume molaire.}
\prerequis{Groupes caractéristiques et familles fonctionnelles.}
\end{ouverture}
\essentiel
\essSection{1}{Absorbance et loi de Beer-Lambert}
\begin{definition}[Absorbance]
Une espèce colorée absorbe une partie de la lumière qui traverse sa solution. Le spectrophotomètre
compare l'intensité lumineuse $I_0$ reçue en l'absence d'espèce absorbante à l'intensité $I$ transmise par la
solution ; l'absorbance $A=\log\left(I_0/I\right)$ est une grandeur positive et sans unité. Le spectre
d'absorption $A=f(\lambda)$ présente un maximum à la longueur d'onde $\lambda_{\max}$.
\end{definition}
\schema{02-spectrophotometre}
\begin{loi}[Loi de Beer-Lambert]
À une longueur d'onde $\lambda$ donnée, l'absorbance d'une solution diluée d'une espèce colorée vaut
\[A=\varepsilon(\lambda)\times\ell\times c,\]
avec $\varepsilon$ le coefficient d'absorption molaire (en L$\cdot$mol$^{-1}\cdot$cm$^{-1}$), $\ell$ l'épaisseur de
solution traversée (en cm) et $c$ la concentration (en \molL). Pour une cuve et une longueur d'onde fixées,
$A=k\,c$. Lorsque plusieurs espèces absorbent, leurs absorbances s'additionnent.
\end{loi}
\begin{attention}
La loi n'est valable que pour des solutions diluées (en pratique, une absorbance inférieure à 1,5 environ),
une lumière monochromatique et une solution limpide, dans laquelle l'espèce dosée ne réagit pas.
Une solution trop concentrée doit être diluée avant la mesure.
\end{attention}
\begin{savoirfaire}[Doser une espèce colorée par étalonnage]
\begin{questions}
\item Tracer le spectre de l'espèce et régler le spectrophotomètre à $\lambda_{\max}$, où la mesure est la plus sensible.
\item Faire le « blanc » avec le solvant, puis mesurer l'absorbance de solutions étalons de concentrations connues.
\item Tracer $A=f(c)$ : les points s'alignent sur une droite passant par l'origine, d'équation $A=k\,c$.
\item Mesurer l'absorbance de la solution à doser, au besoin diluée pour rester dans la gamme d'étalonnage,
et en déduire $c=A/k$.
\end{questions}
\end{savoirfaire}
\schema{02-spectre-etalonnage}
\essSection{2}{Conductivité et loi de Kohlrausch}
\begin{definition}[Conductance et conductivité]
Une solution ionique conduit le courant électrique. La conductance $G=I/U$ (en siemens, S) d'une portion de
solution comprise entre deux plaques dépend de leur aire $S$ et de leur distance $L$ : $G=\sigma\,S/L$.
La conductivité $\sigma$ (en S$\cdot$m$^{-1}$) ne dépend que de la solution et de sa température. Un conductimètre
étalonné affiche directement $\sigma$, souvent en mS$\cdot$cm$^{-1}$.
\end{definition}
\begin{loi}[Loi de Kohlrausch]
La conductivité d'une solution diluée est la somme des contributions de ses ions :
\[\sigma=\sum_i\lambda_i\,[X_i],\]
où $\lambda_i$ est la conductivité molaire ionique de l'ion $X_i$ (en S$\cdot$m$^2\cdot$mol$^{-1}$) et $[X_i]$ sa
concentration (en mol$\cdot$m$^{-3}$). Pour une solution d'un seul soluté de concentration $c$, $\sigma=k\,c$ :
la conductivité est proportionnelle à la concentration.
\end{loi}
\begin{center}\small\renewcommand{\arraystretch}{1.25}
\begin{tabular}{@{}l*{7}{c}@{}}\toprule
Ion & \ce{H3O+} & \ce{HO-} & \ce{Na+} & \ce{K+} & \ce{Ag+} & \ce{Ca^2+} & \ce{Mg^2+}\\
$\lambda$ (mS$\cdot$m$^2\cdot$mol$^{-1}$) & 35,0 & 19,9 & 5,01 & 7,35 & 6,19 & 11,9 & 10,6\\\midrule
Ion & \ce{Cl-} & \ce{NO3-} & \ce{CH3COO-} & \ce{HCO3-} & \ce{SO4^2-} & \ce{NH4+} &\\
$\lambda$ (mS$\cdot$m$^2\cdot$mol$^{-1}$) & 7,63 & 7,14 & 4,09 & 4,45 & 16,0 & 7,35 &\\\bottomrule
\end{tabular}\par\smallskip
\legende{Conductivités molaires ioniques à 25 °C. Les ions oxonium et hydroxyde conduisent nettement mieux que les autres.}
\end{center}
\begin{attention}
Dans la loi de Kohlrausch, les concentrations s'expriment en mol$\cdot$m$^{-3}$ : 1 mmol$\cdot$L$^{-1}$ = 1 mol$\cdot$m$^{-3}$.
De plus, 1 mS$\cdot$cm$^{-1}$ = 0,1 S$\cdot$m$^{-1}$. La loi n'est valable que pour des solutions diluées,
de concentration inférieure à $10^{-2}$ \molL\ environ.
\end{attention}
\schema{02-cellule}
\essSection{3}{Quantité de matière d'un gaz}
\begin{loi}[Équation d'état du gaz parfait]
\[pV=nRT,\]
avec $p$ la pression (en Pa), $V$ le volume (en m$^3$), $n$ la quantité de matière (en mol),
$T$ la température absolue (en K, $T=\theta+273{,}15$) et $R=\qty{8,314}{J.mol^{-1}.K^{-1}}$.
Le volume molaire $V_\mathrm{m}=RT/p$ ne dépend pas de la nature du gaz : environ \qty{24,0}{L.mol^{-1}} à 20 °C sous 1 bar.
\end{loi}
\begin{attention}
Le modèle du gaz parfait décrit bien les gaz sous faible pression (jusqu'à quelques bars), loin de leur
liquéfaction. Il devient approximatif sous forte pression, par exemple dans une bouteille de plongée.
Convertir systématiquement : 1 bar = $10^5$ Pa, 1 hPa = 100 Pa, 1 L = $10^{-3}$ m$^3$.
\end{attention}
\schema{02-capteur-pression}
\essSection{4}{Spectroscopies UV-visible et infrarouge}
\begin{propriete}[Spectre UV-visible et couleur]
Une solution absorbe surtout autour de $\lambda_{\max}$ ; sa couleur perçue est la couleur complémentaire de la
couleur absorbée. La position de $\lambda_{\max}$ et l'allure du spectre, comparées à des données tabulées,
permettent d'identifier une espèce chimique.
\end{propriete}
\schema{02-cercle}
\begin{definition}[Spectre infrarouge]
Un spectre IR représente la transmittance $T$ (en \%) en fonction du nombre d'onde $\sigma=1/\lambda$, exprimé en
cm$^{-1}$, de 4\,000 à 500 cm$^{-1}$ environ, l'axe étant orienté vers la gauche. Chaque type de liaison absorbe
dans un domaine caractéristique : une bande d'absorption est un « creux » du spectre. Au-dessus de 1\,500 cm$^{-1}$,
les bandes identifient les groupes caractéristiques ; en dessous, la zone d'empreinte est propre à chaque molécule.
\end{definition}
\begin{center}\small\renewcommand{\arraystretch}{1.2}
\begin{tabularx}{\linewidth}{@{}lcY@{}}\toprule
\otTete{Liaison} & \otTete{Nombre d'onde (cm$^{-1}$)} & \otTete{Allure de la bande}\\\midrule
O--H d'un alcool libre & 3\,580 -- 3\,670 & fine, forte\\
O--H d'un alcool lié (liaisons hydrogène) & 3\,200 -- 3\,400 & large, forte\\
O--H d'un acide carboxylique & 2\,500 -- 3\,200 & très large, moyenne\\
N--H d'une amine ou d'un amide & 3\,100 -- 3\,500 & moyenne ; deux bandes pour \ce{-NH2}\\
C--H & 2\,800 -- 3\,100 & moyenne à forte\\
C--H d'un aldéhyde & 2\,650 -- 2\,800 & deux bandes fines, moyennes\\
C=O d'un ester & 1\,700 -- 1\,740 & fine, forte\\
C=O d'un aldéhyde ou d'une cétone & 1\,650 -- 1\,730 & fine, forte\\
C=O d'un acide carboxylique & 1\,680 -- 1\,710 & forte\\
C=O d'un amide & 1\,650 -- 1\,700 & forte\\
C=C & 1\,625 -- 1\,685 & moyenne\\
C--O & 1\,050 -- 1\,450 & forte\\\bottomrule
\end{tabularx}\end{center}
\schema{02-spectres-ir}
\begin{savoirfaire}[Exploiter un spectre infrarouge]
Repérer les bandes situées au-dessus de 1\,500 cm$^{-1}$ et les associer aux liaisons à l'aide de la table.
Utiliser aussi l'absence d'une bande : sans bande vers 1\,700 cm$^{-1}$, la molécule ne possède pas de liaison C=O.
Conclure sur la famille fonctionnelle, puis confronter le résultat aux autres données (formule brute, réaction).
\end{savoirfaire}
\begin{rediger}
Écrire « on observe une bande large et forte vers 3\,300 cm$^{-1}$, caractéristique de la liaison O--H d'un alcool
lié », en précisant toujours la grandeur (nombre d'onde, et non longueur d'onde) et l'unité. Pour un dosage,
justifier le choix de la méthode et vérifier que la mesure se situe dans le domaine de validité de la loi utilisée.
\end{rediger}
\exercices

\rubrique{Découverte}
\begin{blocLivre}{ex-2-1}
\begin{exercice}[1]{Appliquer la loi de Beer-Lambert}
Pour une solution de sulfate de cuivre(II), le coefficient d'absorption molaire des ions \ce{Cu^2+} vaut
$\varepsilon=\qty{12}{L.mol^{-1}.cm^{-1}}$ à 800 nm. La cuve a une épaisseur $\ell=\qty{1,00}{cm}$.
\begin{questions}
\item Calculer l'absorbance d'une solution de concentration $c=\qty{5,0e-2}{mol.L^{-1}}$.
\item Une autre solution a une absorbance de 0,42. Calculer sa concentration.
\item On mesure l'absorbance de la première solution dans une cuve de 2,00 cm. Quelle valeur obtient-on ?
\item Pourquoi ne peut-on pas utiliser directement la loi pour une solution d'absorbance égale à 3,2 ?
\end{questions}
\end{exercice}
\end{blocLivre}
\begin{blocLivre}{ex-2-2}
\begin{exercice}[1]{Calculer une conductivité}
\begin{donnees}
$\lambda(\ce{Na+})=5{,}01$ ; $\lambda(\ce{Cl-})=7{,}63$ ; $\lambda(\ce{K+})=7{,}35$ ; $\lambda(\ce{NO3-})=7{,}14$,
en mS$\cdot$m$^2\cdot$mol$^{-1}$.
\end{donnees}
\begin{questions}
\item Exprimer la concentration $c=\qty{2,0}{mmol.L^{-1}}$ en mol$\cdot$m$^{-3}$.
\item Calculer la conductivité d'une solution de chlorure de sodium de concentration $c$, en mS$\cdot$m$^{-1}$,
puis en mS$\cdot$cm$^{-1}$.
\item À la même concentration, une solution de nitrate de potassium conduit-elle mieux ou moins bien ? Justifier.
\end{questions}
\end{exercice}
\end{blocLivre}
\begin{blocLivre}{ex-2-3}
\begin{exercice}[1]{La quantité de matière d'un gaz}
\begin{donnees}
$R=\qty{8,314}{J.mol^{-1}.K^{-1}}$ ; $M(\mathrm{He})=\qty{4,0}{g.mol^{-1}}$.
\end{donnees}
\begin{questions}
\item Un ballon de 4,2 L contient de l'hélium à 20 °C sous une pression de \qty{1,05e5}{Pa}. Calculer la quantité
de matière d'hélium, puis sa masse.
\item Une bouteille de plongée de 12 L contient de l'air sous 200 bar à 20 °C. Calculer la quantité de matière
donnée par le modèle du gaz parfait. Pourquoi ce résultat n'est-il qu'approché ?
\end{questions}
\end{exercice}
\end{blocLivre}
\begin{blocLivre}{ex-2-4}
\begin{exercice}[1]{De la couleur au spectre}
\begin{donnees}
Couleurs complémentaires : rouge et cyan ; vert et magenta ; bleu et jaune.
Domaines approximatifs : bleu 430 -- 490 nm ; vert 500 -- 560 nm ; jaune 560 -- 590 nm ; rouge 620 -- 750 nm.
\end{donnees}
\begin{questions}
\item Une solution de permanganate de potassium absorbe surtout à 525 nm. Quelle couleur absorbe-t-elle ?
Quelle est sa couleur perçue ?
\item Une solution de tartrazine, colorant alimentaire jaune, absorbe-t-elle plutôt vers 430 nm ou vers 650 nm ? Justifier.
\item À quelle longueur d'onde faut-il régler le spectrophotomètre pour doser chacune de ces espèces ?
\end{questions}
\end{exercice}
\end{blocLivre}
\begin{blocLivre}{ex-2-5}
\begin{exercice}[1]{Lire un spectre infrarouge}
L'espèce X est l'éthanol, l'éthanal ou l'acide éthanoïque. On dispose de son spectre infrarouge.
\schema{02-ex05-ir}
\begin{questions}
\item Associer les bandes observées au-dessus de 1\,500 cm$^{-1}$ aux liaisons correspondantes, à l'aide de la table du cours.
\item Quelle bande, absente du spectre, permet d'éliminer deux des trois espèces ? Identifier X.
\end{questions}
\end{exercice}
\end{blocLivre}
\rubrique{Entraînement}
\begin{blocLivre}{ex-2-6}
\begin{exercice}[2]{Doser un colorant alimentaire}
Une boisson pour sportifs est colorée par le bleu brillant FCF (E133), dont le maximum d'absorption se situe
à 630 nm. On mesure l'absorbance de solutions étalons dans une cuve de 1,00 cm.
\begin{center}\small\renewcommand{\arraystretch}{1.2}
\begin{tabular}{@{}l*{5}{c}@{}}\toprule
$c$ ($\upmu$mol$\cdot$L$^{-1}$) & 1,0 & 2,0 & 4,0 & 6,0 & 8,0\\
$A$ à 630 nm & 0,13 & 0,26 & 0,52 & 0,79 & 1,04\\\bottomrule
\end{tabular}\end{center}
La boisson, diluée deux fois, a une absorbance de 0,62.
\begin{donnees}
$M(\text{E133})=\qty{792,8}{g.mol^{-1}}$ ; dose journalière admissible (DJA) : 6 mg par kilogramme de masse
corporelle et par jour (Autorité européenne de sécurité des aliments, 2010).
\end{donnees}
\begin{questions}
\item Justifier le choix de la longueur d'onde de travail.
\item Montrer que les mesures vérifient la loi de Beer-Lambert et déterminer le coefficient $k$ tel que $A=k\,c$.
\item En déduire la concentration en quantité de matière, puis en masse, du colorant dans la boisson.
\item Quel volume de boisson un adolescent de 60 kg devrait-il boire chaque jour pour atteindre la DJA ? Conclure.
\end{questions}
\end{exercice}
\end{blocLivre}
\begin{blocLivre}{ex-2-7}
\begin{exercice}[2]{Contrôler un sérum physiologique}
Un sérum physiologique est une solution de chlorure de sodium annoncée à 0,9 \% (9,0 g par litre).
On le dilue vingt fois, puis on mesure la conductivité de la solution diluée : $\sigma=\qty{0,928}{mS.cm^{-1}}$.
Une gamme de solutions étalons de chlorure de sodium donne, à la même température :
\begin{center}\small\renewcommand{\arraystretch}{1.2}
\begin{tabular}{@{}l*{5}{c}@{}}\toprule
$c$ (mmol$\cdot$L$^{-1}$) & 2,0 & 4,0 & 6,0 & 8,0 & 10,0\\
$\sigma$ (mS$\cdot$cm$^{-1}$) & 0,246 & 0,488 & 0,727 & 0,962 & 1,195\\\bottomrule
\end{tabular}\end{center}
\begin{donnees}
$M(\ce{NaCl})=\qty{58,44}{g.mol^{-1}}$ ; matériel : pipettes jaugées de 5,0 et 10,0 mL, fioles jaugées de 100,0 et 200,0 mL.
\end{donnees}
\begin{questions}
\item Proposer un protocole de dilution. Pourquoi faut-il diluer le sérum avant la mesure ?
\item Modéliser la gamme par une relation $\sigma=k\,c$ et déterminer $k$.
\item Déterminer la concentration du sérum, puis sa concentration en masse. L'indication de l'étiquette est-elle vérifiée ?
\end{questions}
\end{exercice}
\end{blocLivre}
\begin{blocLivre}{ex-2-8}
\begin{exercice}[2]{Deux colorants dans un sirop}
Un sirop de menthe contient deux colorants : la tartrazine E102 (jaune) et le bleu patenté V E131 (bleu).
On dilue le sirop dix fois et on mesure, dans une cuve de 1,00 cm : $A_{638}=0{,}53$ et $A_{427}=0{,}51$.
\schema{02-ex08-spectres}
\begin{donnees}
Coefficients d'absorption molaire (L$\cdot$mol$^{-1}\cdot$cm$^{-1}$) : tartrazine : $2{,}70\times10^4$ à 427 nm,
négligeable à 638 nm ; bleu patenté : $1{,}06\times10^5$ à 638 nm et $3{,}0\times10^3$ à 427 nm.
\end{donnees}
\begin{questions}
\item Expliquer pourquoi le sirop paraît vert.
\item Pourquoi exploite-t-on d'abord la mesure à 638 nm ? En déduire la concentration du bleu patenté dans le sirop dilué.
\item À l'aide de l'additivité des absorbances, déterminer la concentration de la tartrazine dans le sirop dilué.
\item Quelle erreur relative commettrait-on en attribuant toute l'absorbance à 427 nm à la tartrazine ?
\item En déduire les concentrations des deux colorants dans le sirop.
\end{questions}
\end{exercice}
\end{blocLivre}
\begin{blocLivre}{ex-2-9}
\begin{exercice}[2]{Trois isomères, trois spectres}
Le propanal, la propanone et le prop-2-én-1-ol ont la même formule brute \ce{C3H6O}.
\schema{02-ex09-ir}
\begin{questions}
\item Écrire la formule semi-développée de chaque espèce et nommer sa famille fonctionnelle.
\item Attribuer chaque spectre à une espèce, en justifiant à l'aide des bandes caractéristiques.
\item Quelle bande permet de distinguer le propanal de la propanone ?
\end{questions}
\end{exercice}
\end{blocLivre}
\begin{blocLivre}{ex-2-10}
\begin{exercice}[2]{Doser un comprimé effervescent par une mesure de pression}
Un comprimé effervescent contient de l'hydrogénocarbonate de sodium \ce{NaHCO3} et de l'acide citrique.
On le place dans un erlenmeyer contenant de l'eau, que l'on ferme aussitôt ; un capteur mesure la pression.
L'erlenmeyer contient $V=\qty{1,10}{L}$ de gaz, à 20 °C. La pression passe de 1\,015 hPa à 1\,352 hPa.
\begin{donnees}
\ce{HCO3- + AH -> A- + CO2(g) + H2O(l)} (AH : acide, en excès) ; $R=\qty{8,314}{J.mol^{-1}.K^{-1}}$ ;
$M(\ce{NaHCO3})=\qty{84,0}{g.mol^{-1}}$ ; étiquette : 1,30 g d'hydrogénocarbonate de sodium par comprimé.
\end{donnees}
\begin{questions}
\item Exprimer, puis calculer, la quantité de dioxyde de carbone formée.
\item En déduire la masse d'hydrogénocarbonate de sodium du comprimé et l'écart relatif à l'indication de l'étiquette.
\item Une partie du dioxyde de carbone reste dissoute dans l'eau. Dans quel sens cela fausse-t-il le résultat ?
\end{questions}
\end{exercice}
\end{blocLivre}
\begin{blocLivre}{ex-2-11}
\begin{exercice}[2]{La conductivité d'une eau minérale}
L'étiquette d'une eau minérale indique sa composition, en mg$\cdot$L$^{-1}$ :
\begin{center}\small\renewcommand{\arraystretch}{1.2}
\begin{tabular}{@{}l*{8}{c}@{}}\toprule
Ion & \ce{Ca^2+} & \ce{Mg^2+} & \ce{Na+} & \ce{K+} & \ce{HCO3-} & \ce{SO4^2-} & \ce{Cl-} & \ce{NO3-}\\
mg$\cdot$L$^{-1}$ & 80 & 26 & 6,5 & 1,0 & 360 & 14 & 10 & 3,8\\
$M$ (g$\cdot$mol$^{-1}$) & 40,08 & 24,31 & 22,99 & 39,10 & 61,02 & 96,06 & 35,45 & 62,00\\\bottomrule
\end{tabular}\end{center}
\begin{questions}
\item Calculer la concentration en quantité de matière de chaque ion, en mmol$\cdot$L$^{-1}$.
\item Vérifier l'électroneutralité de l'eau à partir de ces concentrations.
\item Calculer la conductivité prévue par la loi de Kohlrausch, avec les valeurs de $\lambda$ du cours.
\item On mesure $\sigma=\qty{0,58}{mS.cm^{-1}}$ à 25 °C. Proposer une explication à l'écart observé.
\end{questions}
\end{exercice}
\end{blocLivre}
\rubrique{Objectif bac}
\begin{blocLivre}{ex-2-12}
\begin{exercice}[3]{Contrôle qualité d'un antiseptique}
\begin{docu}{1}{Composition d'une solution antiseptique}
Pour 100 mL : solution d'hypochlorite de sodium correspondant à 500 mg de chlore actif ; permanganate de
potassium \ce{KMnO4} : 0,0010 g ; excipients. Le permanganate, qui colore la solution en rose, protège
l'hypochlorite de la lumière. Les teneurs réelles doivent rester à $\pm5$ \% des valeurs annoncées.
\end{docu}
\begin{docu}{2}{Gamme d'étalonnage à 525 nm}
\centering\small\renewcommand{\arraystretch}{1.15}
\begin{tabular}{@{}l*{5}{c}@{}}
$c(\ce{MnO4-})$ ($\upmu$mol$\cdot$L$^{-1}$) & 20 & 40 & 60 & 80 & 100\\
$A$ & 0,049 & 0,095 & 0,145 & 0,191 & 0,241
\end{tabular}\par\raggedright
La solution antiseptique, placée dans la même cuve, a une absorbance $A=0{,}150$ à 525 nm.
\end{docu}
\schema{02-ex12-spectres}
\begin{donnees}
$M(\mathrm{K})=\qty{39,10}{g.mol^{-1}}$ ; $M(\mathrm{Mn})=\qty{54,94}{g.mol^{-1}}$ ; $M(\mathrm{O})=\qty{16,00}{g.mol^{-1}}$.
\end{donnees}
\begin{questions}
\item Justifier le choix de la longueur d'onde de 525 nm et montrer que l'ion hypochlorite ne gêne pas le dosage.
\item Vérifier que la gamme d'étalonnage suit la loi de Beer-Lambert et déterminer $k$.
\item \initiative La solution antiseptique est-elle conforme pour sa teneur en permanganate de potassium ?
\end{questions}
\end{exercice}
\end{blocLivre}
\begin{blocLivre}{ex-2-13}
\begin{exercice}[3]{Le gonflage d'un coussin de sécurité}
\begin{docu}{1}{Un gonflage en quelques dizaines de millisecondes}
Lors d'un choc, un capteur déclenche la décomposition d'un solide qui libère du diazote. Dans les premiers
générateurs, ce solide était l'azoture de sodium : \ce{2 NaN3(s) -> 2 Na(s) + 3 N2(g)}. Le sodium formé est
ensuite neutralisé par d'autres réactifs. L'azoture, toxique, a depuis été remplacé dans la plupart des générateurs.
\end{docu}
\begin{docu}{2}{Caractéristiques retenues}
Volume du coussin gonflé : 60 L. Juste après le gonflage, on admet que le diazote est à 80 °C sous une
pression de \qty{1,20e5}{Pa}.
\end{docu}
\begin{donnees}
$R=\qty{8,314}{J.mol^{-1}.K^{-1}}$ ; $M(\mathrm{Na})=\qty{22,99}{g.mol^{-1}}$ ; $M(\mathrm{N})=\qty{14,01}{g.mol^{-1}}$.
\end{donnees}
\begin{questions}
\item Le modèle du gaz parfait est-il adapté pour décrire le diazote du coussin ? Justifier.
\item Calculer la quantité de diazote contenue dans le coussin gonflé.
\item \initiative Déterminer la masse d'azoture de sodium que devait contenir le générateur.
\item Après le choc, le diazote se refroidit à 20 °C. Quelle serait sa pression si le volume restait de 60 L ?
Pourquoi le coussin se dégonfle-t-il alors par ses évents ?
\end{questions}
\end{exercice}
\end{blocLivre}
\begin{blocLivre}{ex-2-14}
\begin{exercice}[3]{De la menthe au laboratoire}
\begin{docu}{1}{Une oxydation contrôlée}
Le menthol, extrait de l'huile essentielle de menthe, peut être oxydé en menthone, une cétone utilisée en parfumerie.
Le produit brut est analysé par spectroscopie infrarouge avant purification.
\end{docu}
\schema{02-ex14-menthol}
\begin{questions}
\item Recopier les formules topologiques. Entourer et nommer le groupe caractéristique de chaque molécule.
\item Écrire la formule brute du menthol et celle de la menthone. La transformation modifie-t-elle la chaîne carbonée ?
\item Attribuer les spectres 1 et 2 au menthol et à la menthone, en justifiant.
\item \initiative Que peut-on conclure du spectre 3 du produit brut ? Faut-il purifier le produit ?
\end{questions}
\end{exercice}
\end{blocLivre}
\begin{blocLivre}{ex-2-15}
\begin{exercice}[3]{La salinité de l'eau de mer}
On cherche à estimer la salinité d'une eau de mer, en la traitant comme une solution de chlorure de sodium.
On la dilue cent fois et on mesure : $\sigma=\qty{0,712}{mS.cm^{-1}}$. On utilise la même gamme d'étalonnage que
dans l'exercice 7, modélisée cette fois par une fonction affine.
\begin{codecote}[0.68]
\begin{pythoncol}
import numpy as np
c = np.array([2.0, 4.0, 6.0, 8.0, 10.0])
s = np.array([0.246, 0.488, 0.727, 0.962, 1.195])
a, b = np.polyfit(c, s, 1)
print(round(a, 4), round(b, 3))
print(round((0.712 - b) / a, 2))
\end{pythoncol}
\cote
\begin{consolecol}
0.1186 0.012
5.9
\end{consolecol}
\end{codecote}
\begin{donnees}
$M(\ce{NaCl})=\qty{58,44}{g.mol^{-1}}$ ; salinité moyenne des océans : environ 35 g de sels par kilogramme d'eau ;
$\lambda(\ce{Na+})=5{,}01$ et $\lambda(\ce{Cl-})=7{,}63$ mS$\cdot$m$^2\cdot$mol$^{-1}$.
\end{donnees}
\begin{questions}
\item \capacite{Capacité numérique}Que représentent les deux nombres affichés à la première ligne ? Et la valeur
affichée à la seconde ligne ?
\item En déduire la concentration « équivalente en chlorure de sodium » de l'eau de mer, puis sa salinité en g$\cdot$L$^{-1}$.
\item La loi de Kohlrausch prévoit quelle conductivité pour l'eau de mer non diluée ? Pourquoi fallait-il diluer ?
\item \initiative Discuter la validité du résultat, sachant que l'eau de mer contient aussi d'autres ions
(\ce{Mg^2+}, \ce{SO4^2-}, \ce{Ca^2+}, \ce{K+}…).
\end{questions}
\end{exercice}
\end{blocLivre}
